% !TEX program = xelatex
% 编译方式：在本目录执行两遍 xelatex（详见 docs/_manual_common/README.md）。
% 源码依据：src/dynamics/ 与 include/hacdcpf/dynamics/（含 devices/ integration/ solvers/）源码及
%           docs/transient_runtime.md，公式与参数以代码为准。
\documentclass[UTF8,fontset=fandol,a4paper,11pt]{ctexart}
\usepackage{../../_manual_common/hysim_manual}

\title{\textbf{交直流混合高弹性能源电力系统仿真平台}\\[0.5em]
       {\Large 机电暂态（动力学）技术手册}\\[0.3em]
       {\large 数学模型、接口参数、验证校核与工业级评价}}
\author{HySim-XJTU-HRPES（西安交通大学 HRPES 团队）}
\date{\today}

\begin{document}
\maketitle
\begin{abstract}
\noindent 本手册依据 HySim-XJTU-HRPES 平台机电暂态（动力学）模块
（\srcpath{src/dynamics/}、\srcpath{include/hacdcpf/dynamics/}，含 \srcpath{devices/}、
\srcpath{integration/}、\srcpath{solvers/} 子目录）整理，命名空间 \texttt{hacdcpf::dynamics}，
总入口头 \srcpath{include/hacdcpf/dynamics/dynamics.hpp}。手册覆盖相量域机电暂态稳定 DAE 的
两类求解架构（分块积分与质量矩阵 DAE）、丰富设备模型库与小信号/模态分析。全部结论以源码与
\srcpath{docs/transient\_runtime.md} 契约为准：暂态为相量域方法，支路 P/Q 仅在求解器记录端口
电流后可用，电压健康失败与初始化不收敛保持为可见错误，不被改写为收敛。版式与《碳流追踪与碳排放
核算技术手册》一致。
\end{abstract}

\tableofcontents
\newpage

\section{阅读指南与约定}
\subsection{数据来源}
\begin{itemize}
  \item \srcpath{include/hacdcpf/dynamics/DynamicSolver.hpp} —— 暂态主入口、事件回放、小信号与
        刚性求解器选择声明；
  \item \srcpath{include/hacdcpf/dynamics/DynamicSolverOptions.hpp} —— 求解器类型、步长、容差、
        雅可比模式等选项；
  \item \srcpath{include/hacdcpf/dynamics/DynamicResults.hpp} —— 结果、快照、模态摘要与孤岛频率结构；
  \item \srcpath{include/hacdcpf/dynamics/integration/}、\srcpath{solvers/} —— 显式/隐式积分器与
        质量矩阵 DAE、稀疏牛顿求解；
  \item \srcpath{include/hacdcpf/dynamics/devices/} —— 同步机、励磁、调速、构网/跟网逆变器、VSC、
        电池、IEEE 1547 保护模型；
  \item \srcpath{docs/transient\_runtime.md} —— 运行时契约与边界；
  \item \srcpath{tests/test\_transient\_dynamics.cpp} 等 —— 验证依据。
\end{itemize}
命名空间为 \texttt{hacdcpf::dynamics}。

\subsection{单位制与索引空间}
时间单位 s、频率 Hz、ROCOF Hz/s、相角 rad、功率 MW（或标幺 pu）、电压相量 pu。状态向量
\fld{state}（\texttt{VectorXd}）按设备状态排布；交流三相电压 \fld{vac\_abc}（\texttt{VectorXcd}）、
直流电压 \fld{vdc} 分列；每个孤岛给出重心频率 \fld{coi\_frequency\_hz} 与 \fld{coi\_rocof\_hz\_s}，
孤岛频率见 \fld{island\_frequencies}。索引空间与动态设备装配顺序一致。

\subsection{诚实口径}
暂态为相量域方法：暂态帧不得用静态潮流公式计算支路 P/Q，动态支路潮流只有在求解器记录端口
电流/功率后才出现（此前 \texttt{capabilities.branch\_power=false}）。与 GridLAB-D/OpenDSS 的
对比仅校核共享的平衡快照范围，并不认证逆变器、电机、直流环节、储能或事件方程。电压健康失败与
初始化不收敛保持为可见错误，不得改写为收敛。

\section{功能定位与架构}
\subsection{公共入口族}
\begin{center}\footnotesize
\begin{tabular}{@{}L{6.6cm}L{3.0cm}L{5.4cm}@{}}
\toprule
入口 & 定义位置 & 说明\\
\midrule
\srcpath{run\_transient\_simulation} & dynamics/DynamicSolver.hpp & 一次调用完成暂态仿真：建系、潮流初始化、时域积分\\
\srcpath{DynamicSolver::solve} & 同上 & 在已建好的 \texttt{DynamicSystem} 上求解\\
\srcpath{applyDynamicEventsThrough} & 同上 & 事件回放：按时刻施加扰动并记录\\
\srcpath{summarize\_small\_signal} & 同上 & 当前工作点的小信号/模态筛查\\
\srcpath{maybe\_auto\_select\_stiff\_solver} & 同上 & 成本排序的稳定积分器阶梯自动选择\\
\srcpath{to\_csv} / \srcpath{to\_modal\_report} & dynamics/DynamicResults.hpp & 时序结果与模态报告渲染\\
\bottomrule
\end{tabular}
\end{center}

\subsection{关键数据结构}
\compmeta{DynamicSolverOptions}{include/hacdcpf/dynamics/DynamicSolverOptions.hpp}
主要字段：\fld{solver\_type}（\texttt{DynamicSolverType}：PartitionedEuler/Heun/RK4、
BackwardEulerNewton、TrapezoidalNewton、RosenbrockEuler、MassMatrixDae）、
\fld{t\_start\_s}/\fld{t\_end\_s}/\fld{dt\_s}、\fld{abs\_tol}/\fld{rel\_tol}/\fld{newton\_tol}、
\fld{dae\_jacobian\_mode}（FiniteDifference/\allowbreak HybridAnalytic/\allowbreak HybridAnalyticColored）、
\fld{auto\_select\_stiff\_solver}、\fld{compute\_small\_signal}、\fld{enable\_der\_protection}。

\compmeta{DynamicResults}{include/hacdcpf/dynamics/DynamicResults.hpp}
主要字段：\fld{success}、\fld{snapshots}（\texttt{DynamicSnapshot} 序列）、\fld{initialization}、
\fld{modal}（\texttt{DynamicModalSummary}）、\fld{applied\_event\_records} 与求解统计
\fld{newton\_iterations}、\fld{linear\_factorizations}、\fld{rejected\_steps}。快照
\texttt{DynamicSnapshot} 含 \fld{time\_s}、\fld{state}、\fld{vac\_abc}、\fld{vdc}、
\fld{coi\_frequency\_hz}、\fld{coi\_rocof\_hz\_s}、\fld{island\_frequencies}、\fld{device\_outputs}；
模态摘要含特征值、\fld{damping\_ratio}、\fld{min\_damping\_ratio}、\fld{stable}；孤岛频率
\texttt{DynamicIslandFrequency} 含各岛重心频率与 \fld{total\_inertia\_mws}。

\section{源码等价的求解与观测模型}
\label{sec:dynamics-source-equivalent}

本章只转写当前源码实际执行的离散方程、矩阵装配与判据。符号不用于替换代码中的分支：
$n_x$、$n_{ac}$、$n_{dc}$ 分别对应 \texttt{DaeLayout::n\_x}、\texttt{n\_ac}、
\texttt{n\_dc}；打包未知量严格按
\begin{equation}
 \mathbf u=[\mathbf x;\Re(\mathbf V_{ac});\Im(\mathbf V_{ac});\mathbf V_{dc}]
 \in\mathbb R^{n_x+2n_{ac}+n_{dc}}
\end{equation}
排列。依据：\srcpath{src/dynamics/DynamicSolver.cpp:DaeStepFormula} 与
\srcpath{src/dynamics/SmallSignal.cpp:DaeLayout}。

\subsection{质量矩阵入口实际求解的残差}
令 $\theta$ 对应 \texttt{DaeStepFormula::theta}，$\mathbf f^{-}$ 对应
\texttt{f\_prev}。在 $t+\Delta t$ 处，每次牛顿迭代重新调用所有设备的
\texttt{stamp} 与 \texttt{computeDerivatives}，然后构造
\begin{align}
 \mathbf R_x &=(\mathbf x-\mathbf x^{-})-\Delta t
 \left[\theta\mathbf f(\mathbf x,\mathbf V,t+\Delta t)
 +(1-\theta)\mathbf f^{-}\right],\\
 \mathbf R_{ac} &=\mathbf I_{ac}-\mathbf Y_{ac}\mathbf V_{ac},\\
 \mathbf R_{dc} &=\mathbf I_{dc}-\mathbf G_{dc}\mathbf V_{dc},\\
 \mathbf R&=[\mathbf R_x;\Re(\mathbf R_{ac});\Im(\mathbf R_{ac});\mathbf R_{dc}].
\end{align}
这就是代码求解对象；本手册不另行引入连续质量矩阵方程。若 $\theta<1$ 且前一步导数为空或
维数不等于 $n_x$，函数直接失败；任一导数或残差非有限也直接失败。依据：
\srcpath{src/dynamics/DynamicSolver.cpp:dae\_compute\_derivatives}。

收敛条件只使用无穷范数：
\begin{equation}
 \|\mathbf R\|_\infty\leq\texttt{newton\_tol}.
\end{equation}
成功时迭代数记为当前零基循环下标加一；超过 \texttt{max\_newton\_iters} 返回“不收敛”。
依据：\srcpath{src/dynamics/DynamicSolver.cpp:dae\_finalize}。

\subsection{雅可比的实际装配}
解析网络块不是抽象的 $\partial\mathbf g/\partial\mathbf V$ 占位符，而是将
$\mathbf Y_{ac}=\mathbf G+\mathrm j\mathbf B$ 按代码写入实数矩阵：
\begin{equation}
 \mathbf J_{ac}=\begin{bmatrix}-\mathbf G&\mathbf B\\-\mathbf B&-\mathbf G\end{bmatrix},
 \qquad \mathbf J_{dc}=-\mathbf G_{dc}.
\end{equation}
微分质量块仅在非纯有限差分模式下写入 $\mathbf I_{n_x}$；设备块来自每个设备的
\texttt{addJacobian}。依据：\srcpath{src/dynamics/DynamicSolver.cpp:dae\_add\_analytic\_network\_block}、
\srcpath{src/dynamics/DynamicSolver.cpp:dae\_build\_and\_factor\_jacobian}。

未被解析块覆盖的列采用前向差分，步长逐列为
\begin{equation}
 h_j=\sqrt{\epsilon_{mach}}\max(1,|u_j|),\qquad
 J_{ij}^{FD}=\frac{R_i(\mathbf u+h_j\mathbf e_j)-R_i(\mathbf u)}{h_j}.
\end{equation}
当启用着色有限差分时，同一组中解析行模式不相交的列同时扰动；代码仍按各列的 $h_j$ 回填其
模式包含的行。解析值与差分值的差小于等于 $10^{-14}$ 时不追加稀疏元。依据：
\srcpath{src/dynamics/DynamicSolver.cpp:add\_fd\_remainder\_value（lambda）}。

缓存雅可比及其 LU 只在以下条件全部成立时复用：选项开启、缓存有效、未强制重建、维数相同、
\begin{align}
 |\theta-\theta_{old}|&\leq10^{-14},\\
 |\Delta t-\Delta t_{old}|&\leq10^{-12}
 \max(1,|\Delta t_{old}|,|\Delta t|),\\
 n_{accepted}&<\max(0,\texttt{dae\_jacobian\_max\_reuse\_steps}).
\end{align}
依据：\srcpath{src/dynamics/DynamicSolver.cpp:dae\_cache\_usable}。

\subsection{牛顿方向、回溯与最小二乘回退}
常规方向由已分解的雅可比求解
\begin{equation}
 \mathbf J\Delta\mathbf u=-\mathbf R.
\end{equation}
回溯从 $\alpha=1$ 开始，每次乘 $0.5$，下限为
$\max(10^{-9},\texttt{newton\_damping\_min})$；仅当
\begin{equation}
 \|\mathbf R(\mathbf u+\alpha\Delta\mathbf u)\|_\infty
 \leq(1-10^{-4}\alpha)\max(\|\mathbf R(\mathbf u)\|_\infty,
 \texttt{newton\_tol})
\end{equation}
时接受。依据：\srcpath{src/dynamics/DynamicSolver.cpp:dae\_try\_line\_search}。

若回溯失败，代码最多十次尝试阻尼最小二乘方向：
\begin{align}
 d&=\max\left(1,\max_i |(\mathbf J^T\mathbf J)_{ii}|\right),\\
 (\mathbf J^T\mathbf J+\lambda d\mathbf I)\Delta\mathbf u&=-\mathbf J^T\mathbf R,
 \qquad \lambda\in\{10^{-8},10^{-7},\ldots,10^1\}.
\end{align}
每个方向仍必须通过上一条完全相同的回溯判据；分解、求解或回溯失败后才将 $\lambda$ 乘十。
依据：\srcpath{src/dynamics/DynamicSolver.cpp:dae\_try\_damped\_least\_squares\_direction}。

\subsection{分块积分器的离散式}
下列公式逐项对应 \texttt{system.x.x} 的赋值，不代表同名方法的其他变体。

\paragraph{显式欧拉。}
\begin{equation}
 \mathbf k_1=\mathbf f(t,\mathbf x_0),\qquad
 \mathbf x_1=\mathbf x_0+\Delta t\mathbf k_1,
 \qquad \texttt{dxdt}=\mathbf k_1.
\end{equation}
随后在 $t+\Delta t$ 调用 \texttt{solveNetwork}。依据：
\srcpath{src/dynamics/integration/ExplicitEuler.cpp:step}。

\paragraph{四阶 Runge--Kutta。}
\begin{align}
 \mathbf k_1&=\mathbf f(t,\mathbf x_0),\\
 \mathbf k_2&=\mathbf f(t+\tfrac12\Delta t,\mathbf x_0+\tfrac12\Delta t\mathbf k_1),\\
 \mathbf k_3&=\mathbf f(t+\tfrac12\Delta t,\mathbf x_0+\tfrac12\Delta t\mathbf k_2),\\
 \mathbf k_4&=\mathbf f(t+\Delta t,\mathbf x_0+\Delta t\mathbf k_3),\\
 \mathbf x_1&=\mathbf x_0+\frac{\Delta t}{6}
 (\mathbf k_1+2\mathbf k_2+2\mathbf k_3+\mathbf k_4),\\
 \texttt{dxdt}&=(\mathbf k_1+2\mathbf k_2+2\mathbf k_3+\mathbf k_4)/6.
\end{align}
依据：\srcpath{src/dynamics/integration/RK4.cpp:step}。

\paragraph{后向欧拉。}
初值为 $\mathbf x=\mathbf x_0+\Delta t\mathbf f_0$，迭代残差和雅可比为
\begin{equation}
 \mathbf r=\mathbf x-\mathbf x_0-\Delta t\mathbf f(t+\Delta t,\mathbf x),
 \qquad \mathbf J_r=\mathbf I-\Delta t\mathbf J_f.
\end{equation}
收敛、回溯系数与接受条件和前述 DAE 回溯完全相同，但状态雅可比由
\texttt{numerical\_state\_jacobian} 提供。依据：
\srcpath{src/dynamics/integration/BackwardEuler.cpp:step}。

\paragraph{梯形法。}
初值同后向欧拉，实际残差和雅可比为
\begin{equation}
 \mathbf r=\mathbf x-\mathbf x_0-	frac12\Delta t(\mathbf f_0+
 \mathbf f(t+\Delta t,\mathbf x)),\qquad
 \mathbf J_r=\mathbf I-	frac12\Delta t\mathbf J_f.
\end{equation}
其回溯流程与后向欧拉一致。依据：
\srcpath{src/dynamics/integration/Trapezoidal.cpp:step}。

\paragraph{源码中的 RosenbrockEuler。}
当前实现只有一次线性求解：
\begin{equation}
 (\mathbf I-\Delta t\mathbf J_f)\Delta\mathbf x=\Delta t\mathbf f_0,
 \qquad \mathbf x_1=\mathbf x_0+\Delta\mathbf x,
 \qquad \texttt{dxdt}=\mathbf f(t+\Delta t,\mathbf x_1).
\end{equation}
本手册不把它扩写为多阶段 Rosenbrock 家族。依据：
\srcpath{src/dynamics/integration/Rosenbrock.cpp:step}。

\subsection{小信号模型的实际线性化}
源码首先构造
\begin{equation}
 \mathbf F(\mathbf u)=
 [\mathbf f;\Re(\mathbf I_{ac}-\mathbf Y_{ac}\mathbf V_{ac});
 \Im(\mathbf I_{ac}-\mathbf Y_{ac}\mathbf V_{ac});
 \mathbf I_{dc}-\mathbf G_{dc}\mathbf V_{dc}],
\end{equation}
再将 $\mathbf J=\partial\mathbf F/\partial\mathbf u$ 分块并消元：
\begin{equation}
 \mathbf A=\mathbf f_x-\mathbf f_y\mathbf g_y^{-1}\mathbf g_x.
\end{equation}
代码通过 SparseLU 解 $\mathbf g_y\mathbf Z=\mathbf g_x$，并执行
$\mathbf A\leftarrow\mathbf f_x-\mathbf f_y\mathbf Z$，不显式形成逆矩阵。分解失败或解非有限即失败。
依据：\srcpath{src/dynamics/SmallSignal.cpp:small\_signal\_analysis}。

即使解析装配成功，微分行仍对被设备触及的列用中心差分覆盖：
\begin{equation}
 h_j=10^{-3}\max(1,|u_j|),\qquad
 \mathbf J_{f,j}=\frac{\mathbf f(\mathbf u+h_j\mathbf e_j)-
 \mathbf f(\mathbf u-h_j\mathbf e_j)}{2h_j}.
\end{equation}
代码这样做是为了跨过摆动方程中 $|\text{power imbalance}|\leq10^{-4}$ 的平衡死区；解析装配失败时，
整个 $\mathbf F$ 使用同一步长中心差分。依据：
\srcpath{src/dynamics/SmallSignal.cpp:small\_signal\_analysis}。

对每个特征值 $\lambda_i$，源码输出
\begin{align}
 f_i&=|\Im\lambda_i|/(2\pi),\\
 \zeta_i&=\begin{cases}-\Re\lambda_i/|\lambda_i|,&|\lambda_i|>10^{-12},\\
 1,&|\lambda_i|\leq10^{-12}\text{ 且 }\Re\lambda_i<0,\\
 -1,&\text{其他情况},\end{cases}\\
 \texttt{oscillatory}_i&=(|\Im\lambda_i|>10^{-6}),\\
 \texttt{stable}&=\bigwedge_i(\Re\lambda_i\leq\texttt{stability\_margin}).
\end{align}
正常参与因子先计算 $p_{ki}=|R_{ki}|\,|(R^{-1})_{ik}|$；若 $R^{-1}$ 含非有限值，则改为
$p_{ki}=|R_{ki}|^2$。两种情况都按 $\sum_kp_{ki}$ 归一化，并按阻尼比从小到大排序模态。
依据：\srcpath{src/dynamics/SmallSignal.cpp:small\_signal\_analysis}。

\subsection{频率报告的代码分支}
只使用投运 AC 支路按母线位置建立连通分量。对设备返回的
\texttt{FrequencyParticipation}，惯量权重是 $w_i=H_iS_i$，锚点回退权重是 $S_i$。
有惯量参与者时，岛内结果为
\begin{equation}
 f_{island}=f_n\frac{\sum_iH_iS_i\omega_i}{\sum_iH_iS_i},\qquad
 \mathrm{ROCOF}_{island}=f_n\frac{\sum_iH_iS_i\dot\omega_i}{\sum_iH_iS_i}.
\end{equation}
若惯量和为零但锚点容量和大于零，将 $H_iS_i$ 全部替换为 $S_i$；两者均无时返回
$f_n$ 和 $0$。没有源的连通分量不进入岛列表。系统级结果使用所有连通分量的相同累计量；标称频率
非正时取 $50$。依据：\srcpath{src/dynamics/DynamicFrequency.cpp:computeFrequencyReport}。

\subsection{设备方程覆盖边界}
设备状态导数的权威实现集中在
\srcpath{src/dynamics/devices/BasicDynamicDevices.cpp}。当前实现入口包括：
\texttt{DynamicLoad}、\texttt{DynamicRLLine}、
\texttt{ThreePhaseDynamicLoad}、\texttt{DCDynamicLoad}、
\texttt{DCVoltageSourceDynamic}、\texttt{SynchronousMachine}、
\texttt{FiveMassShaft}、\texttt{Governor}、\texttt{Exciter}、
\texttt{PowerSystemStabilizer}、\texttt{GridFormingInverter}、
\texttt{GridFollowingInverter}、\texttt{VSCConverterDynamic}、
\texttt{PeriodicVariableSourceDynamic}、\texttt{CSVGN1Dynamic}、
\texttt{DERAADynamic}、\texttt{REGCADynamic}、
\texttt{InductionMachineDynamic}、\texttt{ActiveConstantPowerLoadDynamic}、
\texttt{FluxInductionMachineDynamic}、\texttt{DCDCConverterDynamic}、
\texttt{BatteryDynamic}、\texttt{PVDynamic} 与
\texttt{ProtectionRelay}。

本章尚未逐分支转写上述每一种设备的全部导数，但\emph{同步发电机组控制链}已在后续三章逐式转写：
同步电机族见 \S\ref{sec:dynamics-device-machines}、励磁 / AVR 族见 \S\ref{sec:dynamics-device-exciters}、
调速器与 PSS 族见 \S\ref{sec:dynamics-device-governors-pss}，其状态导数、控制框图与参数表均以对应
\texttt{computeDerivatives} 为准。其余设备（跟网 / 构网逆变器、VSC、感应电机、各类负荷、直流环节、
储能、保护继电器等）尚未逐式覆盖，仍以所列 \texttt{computeDerivatives}、对应 \texttt{stamp} 与
\texttt{addJacobian} 为准；此边界用于阻止文档公式被误当作代码不存在的模型。


\section{接口与参数}
\begin{paramtable}
\fld{solver\_type} & — & — & 见正文枚举 & — & 时域积分器（分块或质量矩阵 DAE）\\
\fld{t\_start\_s} & $t_0$ & s & — & — & 仿真起始时刻（另有 \fld{t\_end\_s}）\\
\fld{dt\_s} & $\Delta t$ & s & $10^{-3}\sim10^{-2}$ & — & 积分步长\\
\fld{abs\_tol} & — & — & 小正数 & — & 绝对误差容差（另有 \fld{rel\_tol}）\\
\fld{newton\_tol} & — & — & 小正数 & — & 隐式步牛顿收敛容差\\
\fld{dae\_jacobian\_mode} & — & — & 枚举 & — & 有限差分/混合解析/着色混合解析\\
\fld{auto\_select\_stiff\_solver} & — & — & 布尔 & — & 刚性时自动切换稳定积分器\\
\fld{compute\_small\_signal} & — & — & 布尔 & — & 是否输出小信号/模态摘要\\
\fld{enable\_der\_protection} & — & — & 布尔 & — & 启用 IEEE 1547 DER 保护\\
\fld{success} & — & — & 布尔 & — & 结果：整体是否成功\\
\fld{min\_damping\_ratio} & $\zeta_{\min}$ & — & — & — & 结果：最小阻尼比（模态摘要）\\
\end{paramtable}
求解统计 \fld{newton\_iterations}、\fld{linear\_factorizations}、\fld{rejected\_steps} 随结果返回，
用于评估刚性与步长自适应表现。

\section{验证与边界局限}
\subsection{验证与校核}
注册测试：\srcpath{test\_transient\_dynamics}（\srcpath{tests/test\_transient\_dynamics.cpp}）、
\srcpath{test\_dynamic\_model\_catalog}（\srcpath{tests/test\_dynamic\_model\_catalog.cpp}），以及
ctest 项 \srcpath{transient\_native\_disturbance\_matrix}（标签
\texttt{dynamics;transient;validation}，来自目标 \srcpath{transient\_validation\_matrix}）。

\subsection{边界与局限（如实标注）}
\begin{itemize}
  \item 仅相量域：暂态帧不得用静态潮流公式算支路 P/Q，动态支路潮流须待求解器记录端口电流/功率后
        才可用（此前 \texttt{capabilities.branch\_power=false}）；
  \item 与 GridLAB-D/OpenDSS 的对比仅校核共享的平衡快照范围，不认证逆变器、电机、直流环节、储能或
        事件方程；
  \item 电压健康失败与初始化不收敛保持为可见错误，不得改写为收敛；
  \item 权威文档 \srcpath{docs/transient\_runtime.md}；历史理论见
        \srcpath{docs/archive/theory/dynamics\_electromechanical\_transient\_design.md}。
\end{itemize}
\section{跨模块工业级技术评价基线}

本章规定本手册所述模块进入工程算例、批量研究或生产服务前必须共同满足的评价基线。
具体模型公式、变量、约束和专用算法以正文相应章节为准；本章不扩大源码已经实现的范围。

\subsection{输入、输出、边界条件与数据结构审计}
输入审计应逐项检查有限性、单位、上下界、枚举值和引用完整性。系统对象必须先按所需层级完成
校验；AC 与 DC 母线采用域限定映射，组件稳定 \texttt{index}、向量位置和临时图索引不得混用。
输出审计必须记录单位、索引空间、模型范围、收敛或可行状态、后端、容差、迭代/节点计数和告警。
空系统、空时域、孤岛、重复 ID、零容量、退化约束与不可用可选后端均属于边界测试集。

\subsection{工程假设与数值稳定性分析}
模型使用的平衡、线性化、稳态、独立性、概率分布、时间离散或网络等值假设必须与应用场景匹配。
数值评价至少包含变量缩放、绝对/相对残差、可行性违反、条件数或枢轴质量、步长/阻尼、迭代停滞
和随机种子复现性。若采用正则化、平滑、回退、松弛或近似，应同时报告触发条件、参数值、对主指标
的敏感性以及是否改变模型口径；不得用“已返回结果”替代收敛或有效性证明。

\subsection{复杂度与资源分析}
令 $n_x$、$n_c$、$n_t$、$n_s$、$|V|$、$|E|$ 分别表示变量、约束、时段、场景、节点和边数。
报告应分别给出建模、求解、后处理的时间与内存。图遍历通常为 $O(|V|+|E|)$；逐时段/逐场景
流程至少线性增长为 $O(n_t n_s)$ 次子问题调用；稠密线性代数最坏可达 $O(n_x^3)$，稀疏方法则应
报告结构非零数、填充率和因子分解峰值内存；MILP/NLP 不给出虚假的多项式最坏界，而报告节点数、
间隙、时限和实测扩展曲线。复杂度结论必须基于固定硬件、固定后端和可复现算例族。

\subsection{异常工况处理}
非法输入应在进入求解器前返回结构化错误；不可行、无界、不收敛、时限、内存不足和后端缺失必须
相互区分。部分结果只有在对应 \fld{ValidityFlags}、\fld{model\_scope}、
\fld{model\_limitations} 或等价字段允许时才可使用。异常路径不得静默丢弃 AC/DC 资产、制造平衡
节点、把 NaN 改写为零或把近似解标为全局最优；系统原始对象应保持可恢复和可重试。

\subsection{与其他模块的接口关系}
接口链路遵循“工程语义模型 $\to$ 校验 $\to$ 投影/装配 $\to$ 求解或分析 $\to$ 结果回投影”。
跨模块传递时保留稳定 ID、单位、符号、时间基准和场景身份；消费潮流、OPF、动态或可靠性结果前，
先检查其收敛、覆盖范围和有效性标志。HTTP/JSON/Python 层只做语义保持适配，不重新解释求解结果。

\subsection{测试与验证建议}
最低验证矩阵包括：手算小例、标准算例、边界/异常输入、随机性质测试、算法或后端交叉校核、
序列化往返、稳定 ID 回投影、AC/DC 同号母线、重复运行确定性和规模扩展测试。数值算法还应加入
残差独立重算、雅可比有限差分或自动微分校核；近似模型应与高保真模型比较并给出误差分布，而非
只报告单个平均值。回归报告必须列明构建配置、算例版本、容差、通过/失败/跳过数及未验证范围。

\subsection{工业级评估指标}
\begin{itemize}
  \item \textbf{正确性}：守恒残差、约束最大违反、解析/基准误差、结果回投影一致性；
  \item \textbf{鲁棒性}：收敛率、不可行识别率、异常分类准确率、扰动敏感性与随机复现率；
  \item \textbf{性能}：端到端时延、吞吐量、峰值内存、并行效率与规模增长斜率；
  \item \textbf{可追溯性}：输入模式版本、稳定 ID 覆盖率、模型范围/限制字段完整率、告警可定位率；
  \item \textbf{工程有效性}：与标准、外部引擎或现场数据的偏差，以及误判/漏判对业务指标的影响。
\end{itemize}
指标应预先给出验收阈值和失败处置；未达到阈值时只能标记为实验性或受限适用。

\subsection{适用范围、局限性与后续改进}
本手册只评价当前源码覆盖的模型、后端和数据结构，不对未建模物理过程、未安装求解器或未执行的
外部验证作保证。后续改进应按“缺口 $\to$ 可量化指标 $\to$ 源码/测试变更 $\to$ 独立验证”闭环，
优先消除会造成错误结论的语义缺口，其次提升数值鲁棒性与性能，最后扩展模型范围。


\end{document}
